\footnotesize
\setlength\LTleft{0pt} \setlength\LTright{0pt}

\vspace*{-1\baselineskip}
\subsection{Siglen und editorische Zeichen}
\vspace{1ex}

\begin{tabular}{lp{110mm}}
\textit{A} 		& fremdhändige Abschrift eines fremden Textes\\
\textit{E}, \textit{E}\textsuperscript{1} 		& Erstdruck\\
\textit{E}\textsuperscript{2} \ldots 		& weitere Drucke\\
\textit{L}, \textit{L\textsuperscript{1}} \ldots & Leibniz, eigenhändig\\
\textit{l} 		& Leibniz, Abschrift von Schreiberhand\\
\textit{LiA} 	& Leibnizens eigenhändige Bemerkungen und Verbesserungen in einer fremdhändigen Abschrift eines fremden Textes\\
\textit{Lil} 		& Leibnizens eigenhändige Bemerkungen und Verbesserungen in einer Abschrift von Schreiberhand\\
\textit{Ü} 		& Übersetzung\\
\lbrack~\rbrack	& bei Datierungen: erschlossenes Datum
			\newline im Text und bei Abbildungen: Änderungen, Ergänzungen und Erläuterungen des Herausgebers
			\newline von Leibniz benutzte eckige Klammern werden im Erläuterungs\-apparat angezeigt\\
\lbrack/\rbrack	& Zeilenfall in diplomatisch wiedergegebenem Text %
			oder im Stufenapparat \\
\lbrack...\rbrack & im Stufenapparat: Gültiger Text, der im Haupttext vollständig wiedergegeben wird\\
\lbrack\textit{!}\rbrack & im Stufenapparat: nicht verbesserte Rechen- oder Sprachfehler\\ %
\textlangle\ \textrangle & Konjektur schwer lesbarer oder durch Beschädigung der Handschrift ausgefallener Wörter bzw. Wortteile\\
\textlangle...\textrangle\ \textlangle\textendash\textrangle\ \textlangle\textendash\ \textendash\textrangle
			& nicht entzifferter bzw. durch Beschädigung der Handschrift ausgefallener Text 					unbestimmter Länge oder im Umfang vermutlich eines oder mehrerer Worte\\
\textit{Kursivierung} 	& zitierter Titel von Büchern oder Schriften
				\newline wörtliches oder fast wörtliches Zitat;
				als \glqq fast wörtlich\grqq\ gilt eine Text\-wiedergabe, die unbedeutsam von der Vorlage 				abweicht, etwa bei flüchtiger Wortfolge oder Kasus\-änderungen\\
\textls{Sperrung}	& Hervorhebungen durch Leibniz %
\end{tabular}

\vspace{1ex}


\subsection{Abkürzungen \lowercase{(allgemein)}}
\vspace{-2.5ex}


\begin{longtable}{ll}
a.a.O. 	& am angegebenen Ort\\
Anm. 	& Anmerkung\\
art. 		& articulum (articula), article(s)\\
Aufl. 		& Auflage\\
Bd(e) 	& Band (Bände)\\	%
bes. 		& besonders\\
Bl. 		& Blatt\\
bspw. 	& beispielsweise\\
bzw.  	& beziehungsweise\\
c., cap.	& caput (capita), capitulum (capitula)\\
ca. 		& circa\\
chap. 	& chapitre(s)\\
def. 		& definitio (definitiones), definition(s)\\
d.h. 		& das heißt\\
ebd. 		& ebenda\\
e.g. 		& exempli gratia\\
eigh.		& eigenhändig\\
erg. 		& ergänzt\\
Erl. 		& Erläuterung\\
f. (ff.) 	& folgend(e)\\
Fig. 		& Figur\\
fl. 		& floruit\\
fol. 		& folio\\
gestr. 	& gestrichen\\
GWLB 	& Gottfried Wilhelm Leibniz Bibliothek \textendash\ Niedersächsische Landesbibliothek\\
H. 		& Hälfte\\
Hrsg. (hrsg.) 	& Herausgeber (herausgegeben)\\
Hs. (Hss.) 	& Handschrift, (Handschriften)\\
Jh. 		& Jahrhundert\\
l. 		& liber (libri), livre(s)\\
LBr 		& \textsc{Hannover}, GWLB, Leibniz-Briefwechsel (mit Signatur)\\
LH 		& \textsc{Hannover}, GWLB, Leibniz-Handschriften (mit Signatur)\\
lib. 		& liber (libri)\\
m. 		& mit\\
Marg. 	& Marginalie(n)\\
Ms. 		& Manuskript; \textsc{Hannover}, GWLB, Handschriften (mit Signatur)\\
n. 		& numerus (numeri)\\
N. 		& Stücknummer(n) in der \textit{LSB}-Ausgabe\\
Nachdr. 	& Nachdruck\\
NB 		& nota bene\\
NLA\;HA 	& Niedersächsisches Landesarchiv, Hannover\\
Nr. 		& Nummer(n)\\
o.S. 		& ohne Seitenangabe\\
p. 		& pagina(e), page(s)\\
r\textsuperscript{o} 	& recto\\
S. 		& Seite(n)\\
Sp. 		& Spalte(n)\\
tlw. 		& teilweise\\
u. 		& und\\
u.a. 		& und andere, unter anderem\\
u.ö. 		& und öfters\\
v. 		& van, von\\
v.c. 		& verbi causa\\
v.g. 		& verbi gratia\\
vgl. 		& vergleiche\\
v\textsuperscript{o} 	& verso\\
Z. 		& Zeile(n)\\
z.B. 		& zum Beispiel
\end{longtable}

\clearpage


\subsection{Abkürzungen (S\lowercase{chriften)}}
\vspace{1ex}

\begingroup		%
\setlength{\leftskip}{10mm}%
\setlength{\parindent}{-\leftskip}
\hspace*{-10mm}%
\textit{AE:} \textit{Acta eruditorum}, hrsg. von O.~Mencke und J.\,B.~Mencke, 50 Bde, Leipzig 1682\textendash1731.\par
\textit{AL1:} \textit{Miscellanea curiosa sive Ephemeridum medico-physicarum Germanicarum academiae Caesareo-Leopoldinae naturae curiosorum ... annus}, 24 Bde, [Wechselnde Verlagsorte] 1670\textendash1706.\par
\textit{AL2:} \textit{Academiae Caesareo-Leopoldinae naturae curiosorum ephemerides, sive Observationum medico-physicarum a celeberrimis viris tum medicis tum aliis eruditis in Germania et extra eam communicatarum centuriae}, 5 Bde, Augsburg/Frankfurt am Main/Leipzig 1712\textendash1722.\par
\textit{AO:} \textsc{Archimedes}, \textit{Opera omnia cum commentariis Eutocii}, hrsg. v.\;J.\,L.\;Heiberg, 3 Bde, Leipzig 1910--1915. Nachdr. Stuttgart 1972.\par
\textit{Berichte:} \textit{Berichte des Secretars der Brandenburgischen Societät der Wissenschaften J.\,Th.\;Jablonski an den Präsidenten G.\,W.\;Leibniz (1700--1715) nebst einigen Antworten von Leibniz}, hrsg. v. A.\;von Harnack, Berlin 1897.\par
\textsc{Bertoloni Meli} 1993: D.\;\textsc{Bertoloni Meli}, \textit{Equivalence and Priority: Newton versus Leibniz. Including Leibniz's Unpublished Manuscripts on the Principia}, Oxford 1993.\par
\textit{BUH:} \textit{Bibliotheque universelle et historique}, hrsg. v.\;J.\;Le Clerc, Amsterdam 1686\textendash1693.\par
\textit{BW:} R.\;\textsc{Boyle}, \textit{The Works}, hrsg. von M.~Hunter und E.\,B. Davis, 14 Bde, London 1999\textendash2000.\par
\textsc{Brather} 1993: H.-S.\;\textsc{Brather}, \textit{Leibniz und seine Akademie. Ausgewählte Quellen zur Geschichte der Berliner Sozietät der Wissenschaften 1697\textendash 1716}, Berlin 1993.\par
Cc 2: \textit{Catalogue critique des manuscrits de Leibniz, Fascicule II (Mars 1672\textendash Novembre 1676)}, hrsg. von A.~Rivaud u.\,a., Poitiers 1914\textendash 1924.\par
\textit{CFP:} G.\,W.\;\textsc{Leibniz}, \textit{Dynamica de potentia et legibus naturae corporeae: Tentamen scientiae novae}, hrsg. von A.~Costa, M.~Fichant und E.~Pasini, 2 Bde, Hildesheim/Zürich/New York 2023.\par
\textit{Chronik:} \textit{Leben und Werk von Gottfried Wilhelm Leibniz: Eine Chronik}, bearb. von K.~Müller und G.~Krönert, Frankfurt am Main 1969.\par
\textit{DL:} R.\;\textsc{Descartes}, \textit{Lettres de M.~Descartes}, hrsg. von C.~de Clerselier, 3 Bde, Paris 1657\textendash 1667.\par
\textit{DO:} R.\;\textsc{Descartes}, \textit{Oeuvres}, hrsg. von C.~Adam und P.~Tannery, 12 Bde, Paris 1879\textendash 1910, 2.~Aufl. ebd.\;1964\textendash 1972.\par
\textsc{Gerland} 1906: G.\,W. \textsc{Leibniz}, \textit{Nachgelassene Schriften physikalischen, mechanischen und technischen Inhalts}, hrsg. von E.~Gerland, Leipzig 1906; Nachdr. Hildesheim und New York 1995.\par
\textit{GO:} G.\;\textsc{Galilei}, \textit{Opere. Edizione Nazionale}, hrsg. von A.~Favaro u.\,a., 20 Bde, Florenz 1890\textendash 1909.\par
\textit{HO:} C.\;\textsc{Huygens}, \textit{Oeuvres complètes}, hrsg. von D.~Bierens de Haan, J.~Bosscha u.\,a., 22 Bde, Den Haag 1888\textendash 1950.\par
\textit{JS:} \textit{Journal des Sçavans}, Paris 1665\,ff.\par
KK 1: \textit{Kritischer Katalog der Leibniz-Handschriften, 1.\;Heft 1646\textendash 1672}, hrsg. von P.~Ritter, als Manuskript veröffentlicht Berlin 1908.\par
\textsc{Lamarra/Palaia} 2005: G.\,W.\;\textsc{Leibniz}, \textit{Essais scientifiques et philosophiques}, hrsg. von A.~Lamarra und R.~Palaia, 3 Bde, Hildesheim, Zürich, New York 2005.\par
\textit{LMG:} \textit{Leibnizens gesammelte Werke}, hrsg. von G.\,H.~Pertz, Dritte Folge: Mathematische Schriften, hrsg. von C.\,I.~Gerhardt, 7 Bde, Berlin, Halle 1849\textendash 1863; Nachdr. Hildesheim 1971.\par
\textit{LOD:} G.\,W.\;\textsc{Leibniz}, \textit{Opera omnia}, hrsg. von L.~Dutens, 6 Bde, Genf 1768; Nachdr. Hildesheim 1990.\par
\textit{LPG:} \textit{Die philosophischen Schriften von Gottfried Wilhelm Leibniz}, hrsg. von C.\,I.~Gerhardt, 7 Bde, Berlin 1875\textendash 1890; Nachdr. Hildesheim 1960\textendash 1961.\par
\textit{LSB:} G.\,W.\;\textsc{Leibniz}, \textit{Sämtliche Schriften und Briefe}, Akademie-Ausgabe, Darmstadt 1923\,ff.
(seit 1954: Berlin).\par
\textit{MB:} \textit{Miscellanea Berolinensia ad incrementum scientiarum, ex scriptis Societatis regiae scientiarum exhibitis edita}, Berlin und Magdeburg 1710--1744.\par
\textit{MOC:} N.\;\textsc{Malebranche}, \textit{Oeuvres Complètes}, hrsg. von A.~Robinet, 23 Bde, Paris 1958\textendash1990.\par
\textit{NRL:} \textit{Nouvelles de la république des lettres}, Amsterdam 1684\textendash 1718.\par
\textit{OF:} \textsc{P.\;de Fermat}, \textit{Œuvres}, hrsg. von P.\;Tannery, Ch.\;Henry. 4 Bände und 1 Supplement. Paris 1891\textendash1912.\par
\textit{PO:} B.\;\textsc{Pascal}, \textit{Œuvres complètes}, hrsg. von J.~Mesnard, 4~Bde, Paris 1964\textendash 1992.\par
\textit{PT:} %
\textit{Philosophical Transactions of the Royal Society}, London 1665\,ff.\par
\textsc{Robinet} 1955: A.\;\textsc{Robinet}, \textit{Malebranche et Leibniz. Relations personnelles}, Paris 1955.\par
\textsc{Robinet} 1988: A.\;\textsc{Robinet}, \textit{G.\,W.\;Leibniz. Iter Italicum (mars 1689 -- mars 1690). La dynamique de la République des Lettres}, Florenz 1988.\par
\textit{SPW}: S.\;\textsc{Stevin}, \textit{The principal works}, hrsg. von E.\,J.\;Dijksterhuis u.\,a., 5~Bde, Amsterdam 1955\textendash1966.\par
\textit{SVF:} \textit{Stoicorum veterum fragmenta}, hrsg. von H.\;von\;Arnim (Bd IV mit M.\;Adler), 4 Bde, Leipzig 1903\textendash1924.\par
\textit{TO:} E.\;\textsc{Torricelli}, \textit{Opere}, hrsg. von G.~Loria und G.~Vassura, 4 Bde, Faenza 1919\textendash 1944.\par
\textit{WO:} J.\;\textsc{Wallis}, \textit{Opera mathematica}, 3 Bde, Oxford 1693\textendash 1699; Nachdr. Hildesheim 1972.\par
\endgroup

\vspace{1ex}


\subsection{Symbole}
\vspace{-2.5ex}

\begin{longtable}{lp{100mm}}
\footnotesize
\Denarius 			& destilletur, distilletur (noch zu bedenken)\\
$\mercury$ 		& Merkur\\
$\rightmoon$ 		& Mond\\
$\odot$ \enskip $\bigodot$ 	& Sonne\\
\end{longtable}



\subsection{Mathematische Zeichen}
\vspace{-2.5ex}

\begin{longtable}{lp{100mm}}
$\smallfrown$ 		& Multiplikation\\[0.5ex]
$\smallsmile$ 		& Division\\[0.5ex]
$\bigtriangledown$ 	& Dreieck\\[0.5ex]%
\textit{f} 			& facit \\[0.5ex]
$\square$ \fbox{2} \enskip qu.	& Quadrat\\[0.75ex]
\fbox{3} \enskip cub. 			& Kubus\\[0.75ex]
$\surd \enskip \sqrt{~~~}$ \enskip Rq. 	& Quadratwurzel\\[0.5ex]
=\ aequ.\ aeq.\ $\sqcap$ $\rightpropto$ $\squaredots$ 	& gleich\\[0.5ex]
$\groesser$ 		& größer als\\[0.5ex]
$\kleiner$ 			& kleiner als\\[0.5ex]
, ,, ,,, $\llcorner~\lrcorner$ 	& Klammerausdrücke\\[0.5ex]
\ovalbox{\makebox[15mm][l]{\phantom{xxx}}} 	& Umrahmungen zur Bezeichnung wegfallender Terme\\[0.5ex]
... ,,, 				& Platzhalter für Terme\\[0.5ex]
$\PM$\quad$\PMG$\quad$\PMH$\quad$\PMD$\quad$\PMB$ & kombinierte Plusminuszeichen%





\end{longtable}

